\documentclass[11pt,a4paper]{article}
\usepackage[margin=1in]{geometry}
\usepackage{amsmath,amssymb}
\usepackage{booktabs}
\usepackage{longtable}
\usepackage{array}
\usepackage{enumitem}
\usepackage{microtype}
\usepackage[hidelinks]{hyperref}
\usepackage{parskip}
\usepackage{titlesec}
\usepackage{fancyvrb}

\titleformat{\section}{\large\bfseries}{\thesection.}{0.6em}{}
\titleformat{\subsection}{\normalsize\bfseries}{\thesubsection}{0.6em}{}
\setlist{itemsep=2pt,parsep=0pt,topsep=4pt}
\renewcommand{\arraystretch}{1.15}

\title{\textbf{Quantitative Anomalies as Market Opportunities}\\[4pt]
\large A Statistical and AI-Assisted Dashboard for Equity and Forex Markets}
\author{Individual Report \\ \small Name: \rule{4cm}{0.4pt} \quad Register No.: \rule{3cm}{0.4pt}}
\date{Submitted up to CIA 2 \\ \small All figures recomputed 26 July 2026}

\begin{document}
\maketitle
\thispagestyle{empty}

\section{Abstract}

The Efficient Market Hypothesis assumes that prices absorb information instantly, that
risk is broadly constant, and that returns cannot be forecast. Fifty years of empirical
work says otherwise, and the documented departures --- lagged macroeconomic transmission,
volatility clustering, fat tails, and cointegration between related assets --- are
collectively called anomalies. This study treats each anomaly not as a defect in the
theory but as a measurable opportunity, and builds a working system that finds them on
live market data.

The project, \textbf{QuantAnomalies}, is a full-stack analytical dashboard with five
statistical pillars: macro factor-and-lag regression, GARCH volatility modelling,
Engle--Granger cointegration for forex pair trading, options pricing under four
closed-form models, and the stochastic processes that underlie all of them. Thirteen
deterministic detectors read these pillars and emit signals; a rule-based decision layer
fuses the signals into a single directional verdict weighted by each detector's own
statistical reliability; a locally hosted large language model then explains the verdict
in plain English without ever being permitted to compute or alter a number.

Validated on the S\&P 500 and a 45-pair forex universe over 2015--2025, the system finds
that standardized macro factors explain 68.7\% of monthly return variation (adjusted
$R^2 = 0.585$), that GARCH(1,1) persistence reaches 0.994 with excess kurtosis of 15.8
and a Jarque--Bera $p$-value indistinguishable from zero, and that USDCHF and USDJPY are
cointegrated at $p = 0.0075$ with a mean-reversion half-life of 69 days. The fused engine
returns a cross-asset scan in 11\,ms warm and a per-asset verdict in 4.5\,ms. The
principal contribution is architectural: statistics detect, deterministic rules decide,
and the language model only explains --- a separation that makes the output identical
whether or not the model is running, and therefore auditable. The acknowledged
limitation, stated deliberately, is that the results are in-sample: no walk-forward
backtest, no multiple-testing correction on the cointegration sweep, and no transaction
costs. These bound what the project currently claims and define the next block of work.

\section{Introduction}

Finance students meet the Efficient Market Hypothesis early and are told, quite
correctly, that it is the right null hypothesis. Its three practical consequences ---
that prices reflect all available information, that volatility is a stable parameter, and
that returns are unforecastable --- are also the three assumptions that empirical finance
has spent decades contradicting. Those contradictions are not scattered noise. They are
repeatable, named, and published: Chen, Roll and Ross showed macroeconomic variables are
priced factors; Mandelbrot, Engle and Bollerslev established that variance moves and
clusters; Engle and Granger showed that individually unforecastable series can be
tethered to each other.

The intellectual move this project makes is small but consequential. An anomaly is a
statement that the model is wrong somewhere. Read from the other direction, it is a
statement about where a careful observer might have an edge. A lagged macro effect is a
positioning window. A volatility regime is a risk-sizing instruction. A stretched
cointegrated spread is a market-neutral bet. The same statistic that a researcher reports
as a rejection of a null is what a practitioner reports as a signal.

Three problems stand between that observation and anything usable. First,
\textbf{fragmentation}: macro-factor work, volatility modelling and cointegration live in
different sub-literatures and different toolchains, so almost nothing lets a user point at
one arbitrary ticker and see all three at once. Second, \textbf{interpretation}: a
persistence of 0.994 or a half-life of 69 days communicates nothing to a non-specialist,
and the statistical literature is not written to fix that. Third,
\textbf{aggregation}: even when several signals are available, combining them is usually
done either by naive averaging, which manufactures false confidence out of disagreement,
or by a machine-learning model whose output is a number nobody can interrogate.

This study addresses all three. It builds an integrated dashboard covering five
statistical pillars on any user-supplied ticker; it adds a natural-language explanation
layer that is constrained to explain rather than compute and is audited against the
figures it was given; and it fuses the signals through an explicit, inspectable rule set
in which conviction \emph{falls} when detectors disagree and dissenting signals are
displayed rather than averaged away.

The scope covered by this report runs from project inception to the CIA 2 presentation:
the data pipeline and its validation, all five analysis pillars, the thirteen-detector
decision layer, the fusion engine, the local AI layer, and the dashboard. Backtesting and
the inferential hardening described in Section~7 are outside this scope and are stated as
open work rather than claimed as done.

\section{Literature Review}

\subsection{Efficient markets and their documented cracks}

Fama (1970, 1991) sets the null: prices fully reflect available information, so
risk-adjusted excess returns are unattainable. The anomalies literature is the record of
where that null fails. Lo and MacKinlay (1988) statistically reject the random-walk model
for U.S. weekly returns. Shiller (1981) shows that prices move far more than dividend
fundamentals can justify --- the excess-volatility finding. These results do not overturn
the hypothesis so much as locate the places where structure survives, and those places are
where this project looks.

\subsection{Macro factors and lagged transmission}

Arbitrage Pricing Theory (Ross, 1976) frames returns as driven by several systematic
factors rather than one market beta. Chen, Roll and Ross (1986), in the canonical
\emph{Economic Forces and the Stock Market}, show that industrial production, inflation
surprises, the term spread and default spreads are priced risk factors for equities. Fama
and French (1989) demonstrate that dividend yield and term/default spreads forecast
returns at business-cycle horizons --- that is, macro effects operate \textbf{with a lag},
which is precisely the effect this project measures. To test direction rather than mere
association, the study uses Granger causality (Granger, 1969), which asks whether a
factor's past improves the forecast of returns beyond what returns' own past already
provides. The macro pillar is therefore an OLS regression of returns on lagged macro
factors, with Granger tests and AIC/BIC lag-depth comparison.

\subsection{Volatility clustering and fat tails}

Mandelbrot (1963), studying cotton prices, observed that large changes are followed by
large changes and small by small, and that returns carry tails far heavier than the normal
distribution admits. Engle (1982) formalised the first half as the ARCH model, letting
conditional variance depend on recent squared shocks; Bollerslev (1986) generalised it to
GARCH, and GARCH(1,1) is the specification fitted here. Cont (2001) catalogues the
stylized facts --- clustering, heavy tails, the leverage effect --- that these models exist
to capture. The properties are tested directly rather than assumed: Ljung--Box (1978) on
squared returns for ARCH effects, and Jarque--Bera (1980) with Q--Q plots for
non-normality. The practical reading is twofold --- a high-volatility regime is a
risk-management and option-pricing signal, and fat tails warn that variance-based risk
measures understate danger.

\subsection{Cointegration and pairs trading}

Engle and Granger (1987) introduced cointegration: two non-stationary price series can
share a stationary linear combination, so they remain tethered in the long run even while
each is individually unpredictable. This is the statistical foundation of statistical
arbitrage. Gatev, Goetzmann and Rouwenhorst (2006), the most-cited empirical treatment,
show that a simple distance-based pairs rule earned significant excess returns on U.S.
equities from 1962 to 2002, and document its decay as the strategy crowded. Vidyamurthy
(2004) provides the standard practitioner pipeline --- hedge ratio, spread, $z$-score,
half-life --- which is the pipeline implemented here.

\subsection{Options pricing and the variance risk premium}

Black and Scholes (1973) and Merton (1973) derived the closed-form European option price
under geometric Brownian motion. Its most useful practical output is implied volatility:
the volatility that equates model price to market price. The gap between option-implied
volatility and subsequently realised or model-forecast volatility --- the variance risk
premium --- is well documented and typically positive (Bakshi and Kapadia, 2003; Carr and
Wu, 2009), meaning options are on average rich. Comparing market implied volatility with
this project's own GARCH forecast is what bridges the volatility pillar to the options
pillar and produces one of the thirteen detector signals.

\subsection{Cross-sectional anomalies}

Basu (1977) found that low price-to-earnings stocks earn higher risk-adjusted returns than
high-P/E stocks --- the value anomaly --- later formalised with the size premium into the
Fama--French three-factor model (1992, 1993) and extended to five factors (2015).
Jegadeesh and Titman (1993) document momentum: 3--12 month winners continue to outperform
over the following 3--12 months, confirmed across asset classes by Asness, Moskowitz and
Pedersen (2013). Thaler (1987) surveys calendar effects. These support the momentum,
relative-strength and seasonality detectors in the decision layer.

\subsection{Statistical and AI tooling}

Econometric work relies on peer-reviewed implementations in \texttt{statsmodels} (Seabold
and Perktold, 2010) and the \texttt{arch} package for GARCH. The novel layer is the use of
a \textbf{local} large language model to translate quantitative output into
natural-language reasoning. The model is deliberately constrained to explain numbers it is
given rather than compute them, mitigating the well-documented hallucination risk in
language generation (Ji et al., 2023) and keeping the analyst note anchored to the
deterministic statistics.

\subsection{Data sources --- assessment}

A study is only as good as its instrument, and here the instrument is a set of
financial-data APIs.

\textbf{FRED}, maintained by the Federal Reserve Bank of St.\ Louis, is an authoritative
and widely cited repository of U.S. macroeconomic series with documented vintages and
revision histories. Its limitation is intrinsic rather than technical: official statistics
are released monthly with genuine publication lags and are subject to later revision,
which is why this study models at monthly frequency and treats macro effects as lagged by
construction.

\textbf{Yahoo Finance}, accessed through \texttt{yfinance}, provides daily OHLCV data
across equities, indices, currencies, futures and crypto, with split- and
dividend-adjusted prices and long history. The documented limitations are occasional data
errors, silent changes to an unofficial API, possible survivorship bias for delisted
names, and adjusted-price conventions that differ between vendors. This project
encountered a more insidious variant, described in Section~5.2, which produced a caveat not
found stated in the sources reviewed: with an unofficial client, the \emph{access pattern}
must be validated, not only the values.

\textbf{DBnomics} re-serves public economic datasets, including FRED, through a single
open API, and is used here purely as a redundancy layer so that one provider outage cannot
halt the study.

\subsection{Research gap}

The anomalies above are individually well studied, but three gaps motivate this work.
\textbf{Integration} --- few accessible tools bring macro-factor, volatility and
cointegration analysis together for an arbitrary user-chosen asset.
\textbf{Interpretation} --- statistical output is opaque to non-specialists, and the remedy
attempted here is a grounded explanation layer that is constrained to explain rather than
compute and is audited against the figures supplied to it. \textbf{Actionability} ---
signals are fused into a single verdict weighted by each detector's own statistical
reliability, with conviction falling on disagreement and dissenting signals displayed
rather than averaged away, so the aggregation stays legible instead of becoming another
opaque score.

\section{Objectives of the Study}

\begin{enumerate}
\item \textbf{Measure macro-driver opportunity.} Identify which macroeconomic forces ---
the volatility index, oil, gold, the dollar index, the 10-year yield, the Fed Funds rate,
inflation and unemployment --- move an asset's returns, and with what time delay, so that
a lag can be read as a positioning window rather than a curiosity.

\item \textbf{Model the risk regime.} Estimate time-varying volatility via GARCH(1,1),
test volatility clustering formally, and quantify the departure from normality, so that
regime changes become an explicit risk-sizing input and tail risk is stated rather than
assumed away.

\item \textbf{Detect mean-reversion opportunity.} Test a 45-pair forex universe for
cointegration, construct the hedge ratio, spread and $z$-score for the strongest pair,
generate entry and exit signals, and estimate the mean-reversion half-life.

\item \textbf{Price options and extract the variance risk premium.} Implement four
closed-form pricing models on a shared engine with full Greeks, wire them to live market
inputs, and compare market implied volatility against the project's own GARCH forecast.

\item \textbf{Synthesise the signals into one decision.} Fuse the detector outputs into a
single directional verdict with explicit conviction, an independent risk axis, and a full
audit trail --- such that conviction falls on genuine conflict instead of averaging
opposing evidence into a confident-looking middle.

\item \textbf{Explain the verdict in natural language, verifiably.} Use a locally hosted
language model to narrate the decision under a strict constraint: it may not compute,
rank, or invent a number, and every figure in its output is checked against the evidence it
was supplied.

\item \textbf{Deliver accessibility and robustness.} Support any user-typed ticker and any
forex combination on a fault-tolerant, cache-validated data pipeline, presented through an
interactive dashboard.
\end{enumerate}

\section{Methodology}

\subsection{Research design}

The design is quantitative, secondary-data, and computational. There is no questionnaire
and no human sample; the instrument is a set of financial-data APIs and the sample is the
live market history they return --- daily price history from 2015 to 2025 across equities,
indices, currencies, futures and crypto, together with monthly U.S. macroeconomic series.
Analysis is conducted at daily frequency for volatility and pairs work and at monthly
frequency for macro regression, matching the publication frequency of official statistics.

\subsection{Data pipeline and validation}

Prices are drawn from Yahoo Finance and macro series from FRED, with DBnomics as a
fallback mirror, and everything is cached as CSV so the system runs offline once populated.
Four validation practices are built in, each prompted by a fault actually encountered:

\begin{itemize}
\item \textbf{Range and sanity validation of every cached series.} This is what first
exposed a corruption in which the volatility and oil caches held each other's data.
\item \textbf{Serialized downloads.} Tracing that corruption revealed the mechanism: the
\texttt{yfinance} client is not thread-safe, and concurrent \texttt{yf.download()} calls
silently return one ticker's data labelled as another's. A deliberately threaded scan
reproduced the fault on demand --- three separate ticker pairs returned byte-identical
frames. The failure is silent and plausible-looking, and no test that merely checks for
missing data would catch it. Downloads now run behind a lock.
\item \textbf{A cache staleness guard.} The download window's end date had been hardcoded,
so once written, a cache was served indefinitely and every ``latest price'' in the
application aged silently --- at discovery, by 208 days. The end date now tracks the
current date, a cached series older than seven days triggers a refetch, and if the refetch
fails the stale cache is still served rather than raising, with the staleness exposed
through a \texttt{data\_freshness()} endpoint so it is visible rather than hidden. Monthly
macro series are exempt, since their lag is genuine.
\item \textbf{Fault-tolerant assembly.} If one factor source fails, that factor is skipped
rather than breaking the entire study, and the omission is recorded.
\end{itemize}

\subsection{Pillar 1 --- Macro factor and lag regression}

A monthly dataset is assembled containing the asset's return alongside eight standardized
macro factors. OLS is estimated with time lags of zero to three months, so that both
contemporaneous and delayed transmission are visible. Because the regressors are
standardized, coefficients are comparable betas and the question ``which driver matters
most'' has a fair answer. Granger-causality tests are run across the factor--lag grid, and
a lagged-correlation heatmap supports visual inspection. Lag depth is compared by AIC and
BIC.

\subsection{Pillar 2 --- GARCH and volatility clustering}

A GARCH(1,1) model is fitted to daily log returns, producing a day-by-day conditional
volatility series and a persistence estimate $(\alpha + \beta)$. Clustering is tested with
the Ljung--Box statistic on squared returns; normality is tested with Jarque--Bera and
inspected with a Q--Q plot alongside skewness and excess kurtosis. The fitted model also
produces the forward volatility forecast consumed by the options pillar.

\subsection{Pillar 3 --- Forex pair trading}

Engle--Granger cointegration is tested across all pair combinations in the forex universe.
For the strongest pair, a hedge ratio is estimated by regression on log prices, the spread
is constructed and standardized to a $z$-score, and entry and exit signals are generated at
thresholds of $\pm 2$ and near zero respectively. The mean-reversion half-life is estimated
from an AR(1) regression on spread changes.

\subsection{Pillar 4 --- Options pricing}

Four closed-form models share a single engine and differ only in the carry term:
Black--Scholes--Merton for equities, Garman--Kohlhagen for FX, Black-76 for futures, and
Bachelier for normal dynamics. First- and second-order Greeks are computed analytically. A
Cox--Ross--Rubinstein lattice handles American early exercise, and two smile-admitting
models are implemented --- the Merton jump-diffusion Poisson series and the Heston
characteristic-function integral. Market wiring uses live spot, a continuously compounded
rate derived from \verb|^TNX|, the actual dividend yield, a GARCH volatility forecast
aggregated over the option's life, and jump parameters fitted from the return
distribution's own tails.

\subsection{Pillar 5 --- Stochastic processes}

Simulation engines are implemented for the Wiener process, geometric Brownian motion,
Ornstein--Uhlenbeck, Cox--Ingersoll--Ross, Merton jump-diffusion and Heston, each reporting
theoretical moments beside empirical ones as a self-check. An Euler-versus-Milstein
convergence study and a variance-reduction comparison (antithetic, control variate, and
both) are included.

\subsection{Coupling between pillars}

The pillars are not independent exercises; four explicit couplings tie them into one
system. The GARCH volatility forecast, aggregated over an option's life, \emph{is} the
sigma fed to Black--Scholes, on the reasoning that an option is a claim on future variance
and a forecast is therefore the correct input rather than a historical average. Heston is
seeded from GARCH by $\kappa = -252\ln(\alpha+\beta)$, with $\theta$ the GARCH long-run
variance and $v_0$ the current conditional variance, since both describe the same dynamics
in discrete and continuous time. The pair spread is fitted as an Ornstein--Uhlenbeck
process whose half-life $\ln(2)/\kappa$ independently reproduces the value the pairs module
derives from AR(1). Finally, market implied volatility minus the GARCH forecast becomes the
variance-risk-premium detector signal.

\subsection{The decision layer}

Three strictly separated layers constitute the design's principal claim.

\textbf{Detection.} Thirteen deterministic detectors emit signals with a severity in
$[0,1]$: trend, breakout, 12-1 momentum and relative performance on the price/trend axis;
mean reversion (RSI plus displacement) and pairs opportunity on the mean-reversion axis;
volatility regime, tail event and options mispricing on the volatility axis; macro
dislocation; and volume anomaly, correlation regime and seasonality on the flow/context
axis.

\textbf{Decision.} The signals are netted into one stance, conviction and position size.
Each signal is weighted by reliability $\times$ severity, where reliability is derived from
the detector's own statistics --- cointegration $p$-value, macro $R^2$, GARCH sample size
--- and never invented. Redundant signals within a family are discounted by
$0.45^{\text{rank}}$, so three views of the same price move do not count as three
confirmations. Direction (agreement $\times$ strength, with strength saturating through
$\tanh$) is separated from size, which is a risk overlay driven by risk-off signals and the
volatility percentile. Conviction is demoted on genuine conflict rather than averaged, and
a full audit trail names every step.

\textbf{Fusion.} A parallel engine answers where the evidence points and how firmly. Eight
of the thirteen detectors feed it. Because detectors speak different direction
vocabularies --- ``uptrend'', ``above\_model'', ``long spread'', ``compressed'', ``rich''
--- a polarity map collapses them onto one bull/bear/neutral axis. Volatility and tail
readings are routed to a separate risk axis so that a volatility spike can never masquerade
as a directional call, and volatility mispricing is kept off both axes since expensive
premium is a relative-value read. With weight $w = \text{severity} \times
\text{reliability}$:

\begin{align*}
\text{tilt}       &= \frac{\text{bull} - \text{bear}}{\text{bull} + \text{bear}} & &\text{direction, } -1..+1\\
\text{agreement}  &= 1 - \frac{\min(\text{bull},\text{bear})}{\max(\text{bull},\text{bear})} & &\text{how one-sided}\\
\text{mass}       &= 1 - e^{-\text{total}/2} & &\text{saturating evidence weight}\\
\text{conviction} &= \text{mass} \times (0.4 + 0.6 \times \text{agreement}) &&
\end{align*}

Strength therefore beats count, nothing pins to 1.0 on noise, and a signal pointing the
other way actively lowers conviction. This replaced an earlier linear form dominated by
signal count that clamped to 1.0 as soon as six of anything appeared.

\textbf{Explanation.} The language model receives the detections and the computed decision
as compact JSON and writes prose. It cannot compute, rank, invent a number, or change the
stance, conviction or size; the system's output is identical whether or not a model is
available. Three mechanisms make that verifiable rather than merely asserted: the prompt
demands a \texttt{STANCE:} line before the prose, so the model's own labelled opinion lands
within roughly ten tokens; that stance renders beside the computed stance, so divergence is
visible rather than hidden; and an \texttt{unverified\_numbers()} check tests every figure
in the narrative against the supplied evidence, surfacing anything unmatched as a grounding
badge in the interface.

\subsection{Implementation and performance}

The backend is FastAPI with pandas, numpy, statsmodels, arch and scipy, exposing 38
endpoints; the frontend is React 19 with Vite, Recharts and framer-motion. The language
runtime is a local \texttt{llama.cpp} server running Qwen3-4B on an RTX 4050 via Vulkan,
with a CPU fallback and a rules-only path if no model is running --- nothing leaves the
machine. Price detectors run broad and cheap, one cointegration sweep serves the whole FX
basket, and GARCH is cached on (ticker, last data date) so it refits once per trading day
regardless of request volume. A detector that fails is recorded in diagnostics rather than
raised, so one dead factor costs one signal rather than the verdict.

\subsection{Verification strategy}

The numerical work is checked rather than trusted, with each check computed at request time
and returned in the API response so it is visible in the interface. Put--call parity is a
model-free arbitrage identity; Greeks are cross-checked against central finite differences;
every Monte Carlo estimator is paired against the exact price for the same model; the
binomial lattice is checked for $1/N$ error decay; and the Heston characteristic function
is checked by driving vol-of-vol to zero, which must collapse the model to Black--Scholes.

\section{Data Analysis and Results (to CIA 2)}

All figures below were recomputed on 26 July 2026 on live data.

\subsection{Pillar 1 --- Macro factor regression (S\&P 500, monthly)}

\begin{center}
\begin{tabular}{lr}
\toprule
\textbf{Quantity} & \textbf{Value} \\
\midrule
$R^2$ & 0.687 \\
Adjusted $R^2$ & 0.585 \\
Granger tests run & 32 \\
Significant Granger relationships & 6 \\
\bottomrule
\end{tabular}
\end{center}

Standardized macro factors explain roughly 69\% of monthly return variation. Because the
regressors are standardized, the coefficients are directly comparable, so the ranking of
drivers is meaningful rather than an artefact of differing units. Six of thirty-two Granger
tests reach significance, indicating that a subset of factors leads returns rather than
merely co-moving with them --- which is the lag that Objective 1 sought to measure. The gap
between $R^2$ and adjusted $R^2$ reflects the cost of the lag structure in degrees of
freedom and is reported rather than suppressed.

\subsection{Pillar 2 --- GARCH and the return distribution (S\&P 500, daily)}

\begin{center}
\begin{tabular}{lr}
\toprule
\textbf{Quantity} & \textbf{Value} \\
\midrule
Sample & 2{,}905 trading days \\
GARCH(1,1) persistence $(\alpha+\beta)$ & 0.994 \\
Excess kurtosis & 15.8 \\
Skewness & $-0.65$ \\
Jarque--Bera $p$-value & $\approx 0$ \\
\bottomrule
\end{tabular}
\end{center}

Persistence of 0.994 means volatility shocks decay very slowly: a turbulent period predicts
further turbulence for weeks, which is exactly the clustering that Mandelbrot described and
that ARCH/GARCH was built to formalise. Excess kurtosis of 15.8 against a normal benchmark
of 0, with negative skew, says that extreme moves --- and particularly extreme
\emph{downward} moves --- occur far more often than a Gaussian assumption admits.
Jarque--Bera rejects normality decisively. The practical consequence is that
variance-based risk measures on this series understate danger, and that a volatility regime
is a legitimate risk-sizing input rather than a descriptive statistic.

One caveat is stated by the code itself and displayed in the interface: as persistence
approaches 1, the GARCH long-run variance becomes ill-conditioned. Short-horizon forecasts
from this fit are reliable; the long-run level is indicative only.

\subsection{Pillar 3 --- Cointegration and pairs (45-pair forex universe)}

\begin{center}
\begin{tabular}{lr}
\toprule
\textbf{Quantity} & \textbf{Value} \\
\midrule
Strongest pair & USDCHF / USDJPY \\
Engle--Granger $p$-value & 0.0075 \\
Hedge ratio & $-0.339$ \\
Mean-reversion half-life & 69 days \\
\bottomrule
\end{tabular}
\end{center}

The pair tested cointegrated on log prices, and the half-life of 69 days sets the horizon
over which a stretched spread would be expected to close. The independent
Ornstein--Uhlenbeck fit in Pillar 5 reproduces this half-life from $\ln(2)/\kappa$, which
is a genuine cross-check: two different estimation routes, one AR(1)-based and one
SDE-based, agree on the same reversion speed.

The result is reported with its qualification. The sweep searches many pair combinations
and reports the best, so the $p$-value is not corrected for multiple comparisons; against a
Bonferroni threshold over the 15 combinations in the reported subset, 0.0075 does not
survive $0.05/15 = 0.00333$. What supports the finding is the economic rationale --- both
legs are USD-based safe-haven crosses and therefore share a common driver --- rather than
the $p$-value in isolation. Benjamini--Hochberg correction is the natural remedy and is
scheduled work.

\subsection{Pillar 4 --- Options and the variance risk premium (AAPL)}

\begin{center}
\begin{tabular}{lr}
\toprule
\textbf{Quantity} & \textbf{Value} \\
\midrule
Market implied volatility & 33.1\% \\
GARCH forecast volatility & 28.8\% \\
Variance risk premium & $+4.3$ points \\
\bottomrule
\end{tabular}
\end{center}

The positive sign matches the documented literature (Bakshi and Kapadia; Carr and Wu):
options are on average rich relative to subsequently realised volatility. Because the
Heston model here is seeded from GARCH rather than calibrated to observed option prices,
the comparison is not circular --- calibrating to prices would have made the premium an
artefact of the fit. The consequence is that the model smile is a prediction rather than a
fit, which is the intended trade-off.

\subsection{Numerical verification}

\begin{center}
\begin{tabular}{p{4.4cm}p{5.6cm}p{4.0cm}}
\toprule
\textbf{Check} & \textbf{Method} & \textbf{Result} \\
\midrule
Put--call parity & model-free arbitrage identity & violation 0.0 \\
Greeks & central finite differences, $O(h^2)$ & max difference $4.6\times10^{-7}$ \\
GBM Monte Carlo & vs analytic Black--Scholes & within 3 standard errors \\
Merton Monte Carlo & vs exact Poisson-weighted series & $\Delta$ 0.053, SE 0.039 \\
Heston Monte Carlo & vs Fourier inversion of the CF & $\Delta$ 0.067, SE 0.050 \\
Heston CF & vol-of-vol $\to 0$ collapses to Black--Scholes & $\Delta\ 3\times10^{-6}$ \\
Binomial lattice & error against $1/N$ & halves per step doubling \\
American call, no dividend & premium must be exactly zero & 0.0 \\
Longstaff--Schwartz & vs 800-step lattice & $\Delta$ 0.023, SE 0.029 \\
Black--Scholes smile & implied-vol spread across strikes & 0.000 \\
\bottomrule
\end{tabular}
\end{center}

The Heston check is the most informative. Driving vol-of-vol to zero must reduce the model
to Black--Scholes, and agreement to six decimal places validates the complex arithmetic,
the branch-cut handling and the quadrature simultaneously --- against a formula derived by
an entirely different route.

Variance-reduction efficiency was measured rather than asserted: plain Monte Carlo
$1.00\times$, antithetic $1.97\times$, control variate $5.01\times$, and both combined
$\mathbf{19.27\times}$.

\subsection{Decision engine performance}

\begin{center}
\begin{tabular}{lr}
\toprule
\textbf{Quantity} & \textbf{Value} \\
\midrule
Detectors implemented & 13 (8 feeding the fusion engine) \\
Cross-asset feed, warm & 11\,ms \\
Per-asset verdict, warm & 4.5\,ms \\
Off-universe asset, cold & 0.9\,s \\
AI narration & $\approx 4$\,s, $\approx 28$ tokens/sec, streaming, local \\
\bottomrule
\end{tabular}
\end{center}

The warm figures are the result of the tiering described in Section~5.10 --- a single
cointegration sweep amortised across the FX basket and a GARCH cache keyed on the last data
date, so the model refits once per trading day rather than once per request.

\subsection{Data-pipeline findings}

Two findings are worth reporting as results in their own right, because both are
methodological rather than incidental. The first is the \texttt{yfinance} thread-safety
fault described in Section~5.2: concurrent downloads silently mislabel one ticker's data as
another's, reproduced on demand with three byte-identical frame pairs. The second is the
frozen end date that had left every cached price series 208 days stale while the dashboard
reported them as current. Both failures were silent and both produced plausible-looking
output, which is precisely why they are worth stating: they were caught by range validation
and a freshness check, not by any test for missing data. A third defect of the same
character --- a dividend-yield unit mismatch applying a 32\% yield to a stock paying 0.32\%
--- was found by the same route.

\subsection{Interpretation against the objectives}

Objectives 1 through 4 are met with quantified results: macro drivers are measured with
their lags and a directional test, the risk regime is modelled and its non-normality
quantified, a cointegrated pair is identified with its reversion horizon, and the variance
risk premium is extracted. Objective 5 is met by the fusion engine, whose behaviour on
conflicting evidence is the specific property that distinguishes it from averaging.
Objective 6 is met and, more importantly, made checkable through the stance-beside-stance
display and the number guardrail. Objective 7 is met: any ticker, any forex combination, on
a pipeline that now validates rather than trusts its sources.

\section{Limitations}

These are stated deliberately, because they bound what the project currently claims.

\begin{enumerate}
\item \textbf{No backtest.} Signals are detected but never evaluated for profit and loss,
Sharpe ratio or drawdown. Everything reported is in-sample. This is the largest gap:
``opportunities'' are not yet evidence.
\item \textbf{Multiple testing is uncorrected.} The cointegration sweep reports the best of
many tests without a Benjamini--Hochberg or Bonferroni correction, as discussed in
Section~6.3.
\item \textbf{OLS standard errors} are used on a monthly regression with autocorrelated
residuals, where Newey--West HAC errors would be appropriate.
\item \textbf{Decision weights are stated judgement, not fitted} --- deliberately, since
fitting them on a few hundred monthly observations would overfit more convincingly than it
would inform.
\item \textbf{No transaction costs} are modelled anywhere.
\item \textbf{GARCH long-run variance is ill-conditioned} at persistence near 1;
short-horizon forecasts hold, the long-run level is indicative.
\item \textbf{A single 10-year rate serves all option maturities}, where a bootstrapped
yield curve would be correct; FX policy rates are approximations.
\item \textbf{Heston is seeded from GARCH, not calibrated to option prices} ---
intentional, to keep the variance-risk-premium comparison non-circular, at the cost of the
smile being a prediction rather than a fit.
\end{enumerate}

\section{Work Planned Beyond CIA 2}

In priority order: a walk-forward backtest of the decision engine, which is the highest
value next step and the one that converts detected signals into evidence;
multiple-comparisons correction on pair selection; Newey--West HAC standard errors on the
macro regression; transaction-cost modelling; a bootstrapped yield curve for option
maturities; and KPSS and structural-break tests to complement ADF.

\section*{References}
\begin{itemize}[leftmargin=1.2em,labelsep=0.4em,label={}]
\item Asness, C. S., Moskowitz, T. J., \& Pedersen, L. H. (2013). Value and momentum everywhere. \emph{Journal of Finance}, 68(3), 929--985.
\item Bakshi, G., \& Kapadia, N. (2003). Delta-hedged gains and the negative market volatility risk premium. \emph{Review of Financial Studies}, 16(2), 527--566.
\item Basu, S. (1977). Investment performance of common stocks in relation to their price-earnings ratios. \emph{Journal of Finance}, 32(3), 663--682.
\item Black, F., \& Scholes, M. (1973). The pricing of options and corporate liabilities. \emph{Journal of Political Economy}, 81(3), 637--654.
\item Bollerslev, T. (1986). Generalized autoregressive conditional heteroskedasticity. \emph{Journal of Econometrics}, 31(3), 307--327.
\item Carr, P., \& Wu, L. (2009). Variance risk premiums. \emph{Review of Financial Studies}, 22(3), 1311--1341.
\item Chen, N.-F., Roll, R., \& Ross, S. A. (1986). Economic forces and the stock market. \emph{Journal of Business}, 59(3), 383--403.
\item Cont, R. (2001). Empirical properties of asset returns: stylized facts and statistical issues. \emph{Quantitative Finance}, 1(2), 223--236.
\item Engle, R. F. (1982). Autoregressive conditional heteroscedasticity with estimates of the variance of United Kingdom inflation. \emph{Econometrica}, 50(4), 987--1007.
\item Engle, R. F., \& Granger, C. W. J. (1987). Co-integration and error correction: representation, estimation, and testing. \emph{Econometrica}, 55(2), 251--276.
\item Fama, E. F. (1970). Efficient capital markets: a review of theory and empirical work. \emph{Journal of Finance}, 25(2), 383--417.
\item Fama, E. F. (1991). Efficient capital markets: II. \emph{Journal of Finance}, 46(5), 1575--1617.
\item Fama, E. F., \& French, K. R. (1989). Business conditions and expected returns on stocks and bonds. \emph{Journal of Financial Economics}, 25(1), 23--49.
\item Fama, E. F., \& French, K. R. (1992). The cross-section of expected stock returns. \emph{Journal of Finance}, 47(2), 427--465.
\item Fama, E. F., \& French, K. R. (1993). Common risk factors in the returns on stocks and bonds. \emph{Journal of Financial Economics}, 33(1), 3--56.
\item Fama, E. F., \& French, K. R. (2015). A five-factor asset pricing model. \emph{Journal of Financial Economics}, 116(1), 1--22.
\item Gatev, E., Goetzmann, W. N., \& Rouwenhorst, K. G. (2006). Pairs trading: performance of a relative-value arbitrage rule. \emph{Review of Financial Studies}, 19(3), 797--827.
\item Granger, C. W. J. (1969). Investigating causal relations by econometric models and cross-spectral methods. \emph{Econometrica}, 37(3), 424--438.
\item Jarque, C. M., \& Bera, A. K. (1980). Efficient tests for normality, homoscedasticity and serial independence of regression residuals. \emph{Economics Letters}, 6(3), 255--259.
\item Jegadeesh, N., \& Titman, S. (1993). Returns to buying winners and selling losers. \emph{Journal of Finance}, 48(1), 65--91.
\item Ji, Z., et al. (2023). Survey of hallucination in natural language generation. \emph{ACM Computing Surveys}, 55(12), 1--38.
\item Ljung, G. M., \& Box, G. E. P. (1978). On a measure of lack of fit in time series models. \emph{Biometrika}, 65(2), 297--303.
\item Lo, A. W., \& MacKinlay, A. C. (1988). Stock market prices do not follow random walks. \emph{Review of Financial Studies}, 1(1), 41--66.
\item Mandelbrot, B. (1963). The variation of certain speculative prices. \emph{Journal of Business}, 36(4), 394--419.
\item Merton, R. C. (1973). Theory of rational option pricing. \emph{Bell Journal of Economics and Management Science}, 4(1), 141--183.
\item Ross, S. A. (1976). The arbitrage theory of capital asset pricing. \emph{Journal of Economic Theory}, 13(3), 341--360.
\item Seabold, S., \& Perktold, J. (2010). statsmodels: econometric and statistical modeling with Python. \emph{Proc. 9th Python in Science Conf.}
\item Shiller, R. J. (1981). Do stock prices move too much to be justified by subsequent changes in dividends? \emph{American Economic Review}, 71(3), 421--436.
\item Thaler, R. H. (1987). Anomalies: the January effect. \emph{Journal of Economic Perspectives}, 1(1), 197--201.
\item Vidyamurthy, G. (2004). \emph{Pairs Trading: Quantitative Methods and Analysis}. Wiley.
\end{itemize}

\end{document}
